% appendix.tex
\chapter{AHP Pairwise Comparison Matrices}
\label{app:ahp}

This appendix documents the pairwise comparison matrices provided by the four
consistent judges (J1, J2, J3, J8) used to derive the final CR-weighted criterion
priorities reported in Table~\ref{tab:weights-main}. All matrices use Saaty's
9-point scale, where entry $a_{ij}$ represents the relative importance of criterion
$i$ over criterion $j$.

\section*{Major Criteria Comparison Matrix}
\label{appsec:major-matrix}

The $4 \times 4$ matrix below shows the aggregated (CR-weighted geometric mean)
pairwise comparisons for the four major criteria: Understandability (C1),
Completeness (C2), Usefulness (C3), and Trustworthiness (C4).

\begin{table}[H]
  \centering
  \caption{Aggregated major-criteria comparison matrix (CR-weighted geometric mean
    of J1, J2, J3, J8).}
  \label{tab:app-major}
  \begin{tabular}{lcccc}
    \toprule
              & C1    & C2    & C3    & C4    \\
    \midrule
    C1 (Under.)  & 1.000 & 0.889 & 1.122 & 0.245 \\
    C2 (Comp.)   & 1.125 & 1.000 & 2.654 & 0.413 \\
    C3 (Use.)    & 0.891 & 0.377 & 1.000 & 0.275 \\
    C4 (Trust.)  & 4.082 & 2.421 & 3.636 & 1.000 \\
    \bottomrule
  \end{tabular}
\end{table}

% ---------------------------------------------------------------
\chapter{Survey Instrument}
\label{app:survey}

\section*{Rating Item Wording}
\label{appsec:items}

Each reviewer was presented with the original MMLU question, the model's selected
answer, and the model's causal explanation. The eight rating items are given below,
each on a 5-point Likert scale (1 = Strongly Disagree; 5 = Strongly Agree).

\begin{enumerate}
  \item \textbf{S1.1 — Clarity:} The explanation clearly describes how the model
    arrived at its answer.
  \item \textbf{S1.2 — Ease of Following:} I can easily follow the reasoning steps
    in the explanation.
  \item \textbf{S2.1 — Coverage:} The explanation addresses all the factors that are
    relevant to the question.
  \item \textbf{S2.2 — Reasoning Depth:} The explanation describes in sufficient
    detail how the result was produced.
  \item \textbf{S3.1 — Goal Support:} The explanation helps me understand the answer
    and achieve the learning goal.
  \item \textbf{S3.2 — Decision-Making / Error Prevention:} The explanation helps me
    judge whether the answer is reliable and to identify potential errors.
  \item \textbf{S4.1 — Reliability:} The explanation is accurate and consistent with
    what I know about this topic.
  \item \textbf{S4.2 — Alignment and Safety:} The explanation clearly indicates when
    the answer is reliable and when it should be treated with caution.
\end{enumerate}

Reviewers also provided an optional free-text comment (maximum 140 characters) and a
confidence slider (0--100\%) for each assessment.

\section*{Attention Check Items}
\label{appsec:attention}

Three attention-check trials were interleaved within each reviewer's task block. Each
check presented a statement with an unambiguously correct rating (e.g., ``This
explanation was written in a language other than English'' when the explanation was in
English). Reviewers who failed two or more attention checks were excluded from
analysis; no reviewers were excluded on this criterion.

% ---------------------------------------------------------------
\chapter{Python Notebooks: Description and Usage}
\label{app:code}

\section*{\texttt{LLMsAnswers.ipynb}}
\label{appsec:llmsanswers}

This notebook loads the MMLU Economics subset using the \texttt{datasets} library,
constructs a structured prompt for each question requesting a JSON-formatted causal
explanation, queries nine local models via the Ollama API, and stores the collected
responses. It also implements the dispersion scoring and MAD-style question selection
described in Section~\ref{sec:meth-qselect}. The notebook produces the set of 12
selected questions and their associated model explanations, which serve as input to
the human evaluation survey.

\textbf{Key dependencies:} \texttt{datasets}, \texttt{subprocess} (Ollama),
\texttt{pandas}, \texttt{sentence-transformers}, \texttt{scikit-learn}.

\section*{\texttt{Priorities\_Calculations.ipynb}}
\label{appsec:priorities}

This notebook implements the AHP consistency analysis reported in
Section~\ref{sec:res-ahp}. It reads the pairwise comparison matrices entered by each
judge, computes the principal eigenvector and eigenvalue, and derives the CR for each
matrix. Matrices satisfying CR $< 0.10$ are retained, and CR-weighted geometric mean
aggregation is applied to produce the final criterion and sub-criterion priorities
reported in Tables~\ref{tab:weights-main}--\ref{tab:weights-sub4}. The notebook
includes a version with the Satisfaction criterion included (for reference) and the
final four-criterion version used in the analysis.

\textbf{Key dependencies:} \texttt{numpy}, \texttt{scipy}, \texttt{pandas}.

\section*{\texttt{E2\_AHP\_CSV\_AND\_GRAPHS.ipynb}}
\label{appsec:e2}

This notebook implements Stage C of the framework. It reads the collected Likert
ratings from \texttt{LikertScale.csv}, computes per-model IQR intervals for each
sub-criterion, applies the AHP weights derived in
\texttt{Priorities\_Calculations.ipynb} to produce aggregated quality intervals $T_m$,
and ranks models by their interval lower bound. It also produces the radar charts
shown in Figures~\ref{fig:radar-qwen14b}--\ref{fig:radar-codellama}, computes
inter-rater reliability statistics (ICC, Krippendorff's $\alpha$), and runs
bootstrapped confidence intervals and effect-size calculations (Cohen's $d$,
Hedges' $g$). Clustering and PCA visualisations of reviewer rating patterns are
also included.

\textbf{Key dependencies:} \texttt{pandas}, \texttt{numpy}, \texttt{matplotlib},
\texttt{scikit-learn}, \texttt{scipy}, \texttt{pingouin} (for ICC).
